\documentclass[11pt,a4paper]{article}
\usepackage[T1]{fontenc}
\usepackage{lmodern,amsmath,amssymb,amsthm,mathtools,microtype}
\usepackage[margin=27mm]{geometry}
\usepackage[colorlinks=true,linkcolor=blue,urlcolor=blue]{hyperref}
\newtheorem{theorem}{Theorem}[section]
\newtheorem{lemma}[theorem]{Lemma}
\newtheorem{corollary}[theorem]{Corollary}
\newtheorem{proposition}[theorem]{Proposition}
\newcommand{\Q}{\mathbb Q}
\DeclareMathOperator{\N}{N}
\title{A single quartic chart for square roots in a quadratic extension}
\author{}\date{8 November 2026}
\begin{document}\maketitle
\begin{abstract}
Let $F$ be a number field and let $K/F$ be quadratic. We construct
a monic quartic with coefficients polynomial in the two components
of an element of $K$, whose nonzero roots in $F$ recover all its
square roots. A cubic change of coordinates ensures that the
chosen rational component never vanishes. The proof uses only
an element of $F$ that is not a cube. It requires neither a choice
among several component charts nor information about rational
points on an auxiliary curve. We obtain a polynomial of total
degree ten, with one quantified variable over $F$, testing whether
a nonzero element of $K$ is a square. Including zero gives two
alternatives: total degree twelve with degree four in the quantified
variable, or total degree ten with degree eight in that variable.
The chart used over $\Q(\sqrt2)$ has integer coefficients given
explicitly; the latter square test has coefficient height eight.
\end{abstract}

\section{A cubic with no nonzero zero}

Let $F$ be a field of characteristic zero, let
$d\in F^\times\setminus F^{\times2}$, and put
$K=F(\rho)$ with $\rho^2=d$. Choose
$\kappa\in F^\times\setminus F^{\times3}$.
For an element $u+v\rho$ of $K$ its \emph{rational component}
means $u\in F$, irrespective of whether $F=\Q$.

\begin{lemma}\label{lem:cubic}
For $z=u+v\rho\ne0$, write
\[
 z^2=U+V\rho,\qquad U=u^2+dv^2,\quad V=2uv.
\]
Define
\begin{equation}\label{eq:multiplier}
 a=U+\frac{\kappa}{2d}V,\qquad
 b=-\frac{\kappa}{d^2}U-\frac12V,\qquad m=a+b\rho.
\end{equation}
Then
\begin{equation}\label{eq:cubic}
 mz=(u^3-\kappa v^3)
       +\left(dv^3-\frac{\kappa}{d^2}u^3\right)\rho.
\end{equation}
In particular the rational component of $mz$ is nonzero.
Moreover $m\ne0$ for every $(U,V)\in F^2\setminus\{(0,0)\}$,
whether or not $U+V\rho$ is a square.
\end{lemma}
\begin{proof}
The rational component of $mz$ is $au+dbv$. Substitution of
$U=u^2+dv^2$ and $V=2uv$ gives $u^3-\kappa v^3$.
The same substitution in $av+bu$ gives the other component
in \eqref{eq:cubic}. If $u^3=\kappa v^3$, either $v\ne0$
and $\kappa=(u/v)^3$, or $v=u=0$. Both are excluded.

The determinant of the linear map $(U,V)\mapsto(a,b)$ is
\[
 \frac{\kappa^2-d^3}{2d^3}.
\]
It is nonzero, since $\kappa^2=d^3$ would make
$d=(\kappa/d)^2$ a square. As $1,\rho$ are linearly independent
over $F$, the multiplier can vanish only when $a=b=0$, and
therefore only when $U=V=0$.
\end{proof}

Every number field admits the required choice of $\kappa$.
Indeed, in its ring of integers choose a nonzero prime ideal
$\mathfrak p$ and an element
$\kappa\in\mathfrak p\setminus\mathfrak p^2$.
Then $v_{\mathfrak p}(\kappa)=1$, whereas the valuation of a cube
is divisible by three. Thus Lemma~\ref{lem:cubic} applies to
every quadratic extension of every number field, with the
coefficient $\kappa$ chosen for that base field.

One may also check the multiplier through its norm:
\begin{equation}\label{eq:norm}
 a^2-db^2=\frac{d^3-\kappa^2}{4d^3}
                              (4U^2-dV^2).
\end{equation}
Both factors on the right are nonzero for $(U,V)\ne(0,0)$.

\section{The quartic and its inverse}

Regard $a,b$ in \eqref{eq:multiplier} as linear polynomials in
$U,V$. Define cubic polynomials over $F$ by
\begin{equation}\label{eq:AB}
 A=(a^2+db^2)U+2dabV,\qquad
 B=(a^2+db^2)V+2abU.
\end{equation}
They satisfy the exact identity
\begin{equation}\label{eq:square}
                  m^2(U+V\rho)=A+B\rho.
\end{equation}
Put
\begin{equation}\label{eq:quartic}
 J(U,V,Y)=Y^4-4A(U,V)Y^2+4dB(U,V)^2.
\end{equation}
This polynomial is monic of degree four in $Y$ and has total
degree at most six.

\begin{theorem}\label{thm:chart}
Let $(U,V)\in F^2\setminus\{(0,0)\}$. Then $U+V\rho$ is
a square in $K$ if and only if there exists $Y\in F^\times$
with $J(U,V,Y)=0$. Every such root recovers a square root by
\begin{equation}\label{eq:inverse}
 z=\frac{a-b\rho}{a^2-db^2}
                  \left(\frac Y2+\frac BY\rho\right).
\end{equation}
Conversely, every square root $z$ of $U+V\rho$ is recovered by
setting $Y$ equal to twice the rational component of $mz$.
Every permitted root $Y$ is simple.
\end{theorem}
\begin{proof}
If $z^2=U+V\rho$, put $z'=mz=u'+v'\rho$.
Lemma~\ref{lem:cubic} gives $u'\ne0$. By
\eqref{eq:square},
\[
       (u')^2+d(v')^2=A,\qquad 2u'v'=B.
\]
Thus $Y=2u'$ is nonzero, $v'=B/Y$, and
$Y^4-4AY^2+4dB^2=0$. The inverse formula follows by
dividing $z'$ by $m$.

Conversely, a nonzero root $Y$ of $J$ satisfies
\[
 \left(\frac Y2+\frac BY\rho\right)^2=A+B\rho,
\]
because the rational component of the square is
$Y^2/4+dB^2/Y^2=A$. The multiplier is nonzero by
Lemma~\ref{lem:cubic}, so \eqref{eq:square} and
\eqref{eq:inverse} give $z^2=U+V\rho$.

At such a root, write $z'=Y/2+(B/Y)\rho=u'+v'\rho$.
Then
\[
 \frac{\partial J}{\partial Y}
    =4Y^3-8AY
    =8Y\bigl((u')^2-d(v')^2\bigr)
    =8Y\N_{K/F}(z')\ne0.
\]
The last inequality follows because $z'=mz$ is a nonzero
element of a field.
\end{proof}

\begin{corollary}\label{cor:test}
For $(U,V)\ne(0,0)$, the element $U+V\rho$ is a square in $K$
if and only if
\begin{equation}\label{eq:reciprocal}
       1-4A(U,V)T^2+4dB(U,V)^2T^4=0
\end{equation}
has a solution $T\in F$. This is a polynomial of total degree
at most ten in $(U,V,T)$. Without excluding zero, the equivalent
test with one quantified variable is
\begin{equation}\label{eq:zero}
 (U^2-dV^2)\bigl(1-4A(U,V)T^2
                         +4dB(U,V)^2T^4\bigr)=0,
\end{equation}
of total degree at most twelve.
\end{corollary}
\begin{proof}
Equation~\eqref{eq:reciprocal} excludes $T=0$, so it is
equivalent to \eqref{eq:quartic} under $Y=1/T$.
Since $d$ is not a square, $U^2-dV^2=0$ holds precisely when
$U=V=0$. This proves \eqref{eq:zero}, including the zero
square.
\end{proof}

\section{A degree-ten test including zero}

If the objective is square membership rather than a finite quartic
chart with rational reconstruction on one open set, a second
formula has lower total degree and smaller coefficients. It does
not require a noncube in the base field.

\begin{proposition}\label{prop:low-height-test}
Let $F$ be any field of characteristic different from two, let
$d\in F^\times$ be nonsquare, and let
$c\in F^\times$ be any nonsquare, for example $c=d$.
An element $U+V\sqrt d$ of $F(\sqrt d)$ is a square, including
zero, if and only if there exists $T\in F$ such that
\begin{equation}\label{eq:low-height-test}
 (4-4UT^2+dV^2T^4)\bigl(V^2-c(U-dT^2)^2\bigr)=0.
\end{equation}
The polynomial has total degree ten and degree eight in $T$.
\end{proposition}
\begin{proof}
A root of the first factor has $T\ne0$ and gives
\[
 \left(\frac1T+\frac{VT}{2}\sqrt d\right)^2=U+V\sqrt d.
\]
Conversely, a square root $u+v\sqrt d$ with $u\ne0$ gives
such a root by taking $T=1/u$.

Since $c$ is nonsquare, the second factor vanishes exactly
when $V=0$ and $U=dT^2$. These are precisely the elements
with a square root of the form $T\sqrt d$, including zero.
Thus the two factors cover all square roots, and each factor
implies square membership. Their product expresses this
alternative using the same quantified variable.
\end{proof}

For $F=\Q$ and $d=2$, choose $c=-1$ and divide the first
factor by two. The test becomes
\begin{equation}\label{eq:height-eight}
        (2-2UT^2+V^2T^4)\bigl(V^2+(U-2T^2)^2\bigr)=0.
\end{equation}
It has twelve monomials and coefficient height eight.
In contrast, the quartic in Theorem~\ref{thm:chart} has
degree four in its finite root and a single rational recovery
formula. These are different useful bounds: the lower-height
test \eqref{eq:height-eight} has degree eight in $T$ and uses
two recovery formulas according to the vanishing factor.

\section{Integer coefficients over $\Q(\sqrt2)$}

Take $F=\Q$, $d=2$, and $\kappa=2$. Multiplying the multiplier
in \eqref{eq:multiplier} by two gives the simpler integral formulas
\begin{align*}
 a&=2U+V, &b&=-(U+V),\\
 A&=6U^3-9UV^2-4V^3,
 &B&=-4U^3+6UV^2+3V^3.
\end{align*}
Thus the explicit monic quartic is
\begin{equation}\label{eq:explicit}
 Y^4-4(6U^3-9UV^2-4V^3)Y^2
                   +8(-4U^3+6UV^2+3V^3)^2=0.
\end{equation}
The nonzero-root condition and inverse formula in
Theorem~\ref{thm:chart} apply with these scaled $a,b,A,B$.
Their elementary identities are
\begin{align*}
 (a+b\sqrt2)(u+v\sqrt2)
      &=2(u^3-2v^3)+(4v^3-u^3)\sqrt2,\\
 a^2-2b^2&=2U^2-V^2,\\
 (a+b\sqrt2)^2(U+V\sqrt2)&=A+B\sqrt2.
\end{align*}
The rational component is nonzero because $2$ is not a rational
cube. For example, a rational cube has $2$-adic valuation
divisible by three, whereas $v_2(2)=1$.

The corresponding one-variable square test, including zero, is
\[
 (U^2-2V^2)\left(1-4(6U^3-9UV^2-4V^3)T^2
                    +8(-4U^3+6UV^2+3V^3)^2T^4\right)=0.
\]
All its coefficients are integers.
The monic quartic \eqref{eq:explicit} and the reciprocal test
for nonzero squares each have ten monomials and coefficient
height $384$. The displayed zero-inclusive test has fifteen
monomials and coefficient height $1056$.

\section{Application to direct elliptic charts}

Put $D=11-6\sqrt2$ and let $E:y^2=x^3-Dx$. Its finite points
with nonzero square abscissa are represented by
\[
                x=h^2,\qquad j^2=h^4-D,\qquad y=hj.
\]
Writing $h=h_0+h_1\sqrt2$ gives
\[
 U=h_0^4+12h_0^2h_1^2+4h_1^4-11,\qquad
 V=4h_0h_1(h_0^2+2h_1^2)+6.
\]
The quantity $h^4-D$ is nonzero: otherwise $h^2=\pm(3-\sqrt2)$;
the negative sign is excluded by the real embeddings, and the
positive sign by the nonsquare norm $7$.
Substitution in \eqref{eq:explicit} therefore represents every
point in this chart by two rational base coordinates $h_0,h_1$
and one finite quartic root. The coefficient degrees in the base
coordinates are at most twenty-four. Both the nonzero-root
condition and its simple-root property are explicit.

The subgroup assertion that these are exactly the nonzero
integer multiples of a specified point is a separate arithmetic
statement. It follows from the stated group hypothesis and the
descent calculation in the companion note
\emph{Square abscissas and direct charts for the elliptic subgroup}.
The present square-root chart and its number-field extension are
unconditional.

\paragraph{Exact verification.}
The accompanying program \texttt{check\_quadratic\_square\_chart.py}
verifies the displayed identities, degrees, monomial counts, and
coefficient heights by exact polynomial arithmetic over $\Q$,
clearing the fixed denominators in $d$ and $\kappa$.
The nonsquare and noncube hypotheses establish nonvanishing in
the proofs above.

\end{document}
